\chapter{Les catégories dérivées.}
\VerdierPhysicalPageAnchor{0174}%
\VerdierSemanticAnchor{ch3d}
\begingroup
\small
\setlength{\parskip}{0pt}
\setlength{\emergencystretch}{1.5em}
\setlength{\abovedisplayskip}{4pt plus 2pt minus 2pt}%
\setlength{\belowdisplayskip}{4pt plus 2pt minus 2pt}%

\VerdierSemanticAnchor{ch3d-1}
\section{Définitions des catégories dérivées.}

\noindent\textit{Dans ce paragraphe, $\mathcal A$ désigne, sauf mention du
contraire, une catégorie abélienne.}

\VerdierSemanticAnchor{ch3d-1-1}
\subsection{Notations et terminologie.}

\VerdierSemanticAnchor{ch3d-1-1-1}
\noindent\textbf{1.1.1.}\quad
Soit $\mathcal A$ une catégorie additive (non nécessairement abélienne). On
désigne par $\operatorname{comp}^{b}(\mathcal A)$ (resp.
$\operatorname{comp}^{+}(\mathcal A)$, resp.
$\operatorname{comp}^{-}(\mathcal A)$) la sous-catégorie pleine de
$\operatorname{comp}(\mathcal A)$ définie par les complexes $X$ tels que
$X^i=0$, pour $|i|$ assez grand (resp. $i$ assez petit, resp. $i$ assez
grand). On désigne par $K^{b}(\mathcal A)$ (resp. $K^{+}(\mathcal A)$,
resp. $K^{-}(\mathcal A)$) la sous-catégorie pleine de $K(\mathcal A)$
définie par les objets de $\operatorname{comp}^{b}(\mathcal A)$ (resp.
$\operatorname{comp}^{+}(\mathcal A)$, resp.
$\operatorname{comp}^{-}(\mathcal A)$).

\VerdierSemanticAnchor{ch3d-1-1-2}
\noindent\textbf{1.1.2.}\quad
Soient maintenant $\mathcal A$ une catégorie (abélienne) et $*$ l'un des
signes typographiques : $b$, $+$, $-$, ``vide''. On désigne par
$\operatorname{comp}^{*,b}(\mathcal A)$ (resp.
$\operatorname{comp}^{*,+}(\mathcal A)$, resp.
$\operatorname{comp}^{*,-}(\mathcal A)$) la sous-catégorie pleine de
$\operatorname{comp}^{*}(\mathcal A)$ définie par les complexes $X$, objets
de $\operatorname{comp}^{*}(\mathcal A)$, tels que $H^i(X)=0$, pour $|i|$
grand (resp.\ pour $i$ petit, resp.\ pour $i$ grand). De même, on désigne par
$K^{*,b}(\mathcal A)$ (resp. $K^{*,+}(\mathcal A)$, resp.
$K^{*,-}(\mathcal A)$) la sous-catégorie pleine de $K^{*}(\mathcal A)$
définie par les complexes $X$, objets de $K^{*}(\mathcal A)$, tels que
$H^i(X)=0$, pour $|i|$ grand (resp. $i$ petit, resp. $i$ grand). On introduit
ainsi huit sous-catégories de $K(\mathcal A)$ :

\VerdierSemanticAnchor{ch3d-1-1-2-diagram}
\begin{center}
  \resizebox{0.88\textwidth}{!}{%
    $\begin{tikzcd}[ampersand replacement=\&,column sep=3.5em,row sep=1.7em]
      \& \& K^{+}(\mathcal A) \arrow[dr] \& \& \\
      \& K^{+,b}(\mathcal A) \arrow[ur] \arrow[dr]
        \& \& K^{,+}(\mathcal A) \arrow[dr] \& \\
      K^{b}(\mathcal A) \arrow[ur] \arrow[dr]
        \& \& K^{,b}(\mathcal A) \arrow[ur] \arrow[dr]
        \& \& K(\mathcal A) \\
      \& K^{-,b}(\mathcal A) \arrow[ur] \arrow[dr]
        \& \& K^{,-}(\mathcal A) \arrow[ur] \& \\
      \& \& K^{-}(\mathcal A) \arrow[ur] \& \&
    \end{tikzcd}$%
  }\hfill\textup{(1.1.2.1)}
\end{center}
\endgroup
